% ID: SEL-D0036-OA-E03
% ESTADO: cotejada. Segunda pasada visual y borrador compilado revisados.
% Fuente: Selectividad/matematicas_aplicadas/Ua1MATEMÁTICAS APLICADAS examen.pdf
% SHA-256: 77e0c295529ccd4ed502b5859d4011cb4b0e955a1e4795e570a9da20ae48293c
% Página: 1; opción A; ejercicio original 3.

\begin{enumerate}[label=\alph*)]
\item \textbf{(1 punto)} Un ilusionista tiene seis cartas: cuatro ases y dos reyes. Saca una carta, la enseña al público y, sin verla, la vuelve a mezclar con las demás. A continuación saca una segunda carta que resulta ser un as. \textquestiondown{}Cuál es la probabilidad de que la primera carta haya sido también un as?
\item \textbf{(1.5 puntos)} Si el ilusionista no devolviera la primera carta a la baraja y la segunda carta extraída fuera un as, \textquestiondown{}cuál es la probabilidad de que la primera carta haya sido también un as?
\end{enumerate}
